%!TEX root = ../thesis-trier.tex

\selectlanguage{ngerman}



Wollten Sie schon einmal wieder von vorne anfangen? Hat sich Ihr Smartphone % mobiltelefon? Handy?
oder Computer schon einmal aufgehängt %aufgehangen  
und nur ein Neustart schuf Abhilfe? Falls ja, so wissen Sie, was es heißt etwas zurückzusetzen, ein zentrales Konzept
der vorliegenden Arbeit.
Im Speziellen handelt sie von der Zurücksetzbarkeit (oder synonym Synchronisierbarkeit) endlicher Zustandsübergangssysteme, auch endliche Automaten genannt.
Automaten interagieren über Eingabefolgen (auch Wörter genannt), welche Zustandsübergänge bewirken, mit ihrer Umwelt
und dienen der Modellierung einer Vielzahl von Systemen wie Schalt\-kreisen, 
Betriebszuständen von Maschinen oder einfachen Agenten in der künstlichen Intelligenz.
%Oft möchte man wissen, ob ein System in einen wohl\-de\-fi\-nierten Zustand versetzbar ist. 
%Beispielsweise könnte ein Schaltkreis nicht mehr synchron mit all seinen Teilen laufen; im Netzwerk könnte der aktuelle Auf\-ent\-halts\-ort unbekannt sein, man möchte aber dennoch die Möglichkeit haben stets einen Weg hinaus zu finden; ein Roboter könnte bei einer Aufgabe festhängen, oder wir kennen einfach nicht seinen aktuellen Zustand, sodass wir ihn in einen wohl\-de\-fi\-nierten Betriebszustand setzen wollen. 
%
%Oft möchte man wissen, ob ein System, beispielsweise einen Schaltkreis dessen Teile
%nicht mehr synchron laufen, in einen wohl\-de\-fi\-nierten Zustand versetzbar ist. 
%Beispielsweise könnte Teile eines Schaltkreises nicht mehr synchron laufen
%oder der aktuelle Zustand einer Maschine ist unbekannt
%und soll zurückgesetzt werden.
%Wir nennen nun einen Automaten synchronisierbar, sofern
Ein Automat heißt synchronisierbar, sofern
es eine Eingabefolge gibt, welche
ihn in einen wohldefinierten Zielzustand versetzt, unabhängig von seinem aktuellen Zustand.
Die Schwierigkeit auf Synchronisierbarkeit zu testen %, sowie die Berechnung einer (möglichst kurzen) synchronisierenden Eingabefolge,
hängt von verschiedenen Parametern ab, zum Beispiel
ob wir das System vollständig kennen, welches Verhalten beobachtbar ist,
ob Nichtdeterminismus erlaubt ist oder ob bestimmte synchronisierende Folgen ausgeschlossen sind.

In dieser Arbeit beschäftigen wir uns mit der Einschränkung 
möglicher synchronisierender Eingabefolgen, das sog.~\emph{constrained synchronization problem} (dt.~\emph{Problem der eingeschränkten Synchronisierbarkeit}).
Es stellt sich nun in natürlicher Weise die Frage, wie die Art der Einschränkung 
die algorithmische Komplexität be\-einflusst. Kann man anhand einer gegebenen Einschränkung
die resultierende Problemstellung in schwierig oder praktikabel lösbar einteilen?
Wir untersuchen komplexitätserhaltende Operationen auf Mengen von erlaubten synchronisierenden Eingabefolgen. Weiterhin
arbeiten wir verschiedene hinreichende Bedingungen  für 
praktikabel lösbare oder schwierige Probleme heraus.
Im Rahmen unserer Untersuchungen tauchen folgende Komplexitätsklassen auf: 
$\AC^0$, \LOGSPACE, \NL, \PTIME, \NP, \PSPACE,
wobei wir Bedingungen für Mitgliedschaft in einer dieser Klassen
oder Schwierigkeit bezüglich \NP\ oder \PSPACE\ 
angeben sowie, außer für~$\AC^0$ und~\PTIME, 
eine Vielzahl von 
vollständigen Problemen für die jeweiligen Komplexitätsklassen.
Die folgenden Arten von Einschränkungen für Mengen synchronisierender Eingabefolgen untersuchen wir: 
endlich, unendlich, spärlich, buchstaben-beschränkt, kommutativ oder durch kleine Automaten beschrieben.
Für kommutative, buchstaben-beschränkte und durch bestimmte kleine Automaton
beschriebene Einschränkungen  
geben wir eine komplette Klassifikation der algorithmischen Komplexität des zugehörigen eingeschränkten Synchronisierbarkeitsproblems. 
Darüber hinaus untersuchen wir die folgenden Klassen von Eingabeautomaten: monotone, partiell geordnete,
kommutative, auflösbare SI- und $\rho$-Automaten.
Unsere Ergebnisse zeigen, dass die resultierenden algorithmischen Komplexitäten
stark vom betrachteten Eingabetyp abhängen. Bei\-spiels\-wei\-se ist das Problem
für kommutative Eingabeautomaten stets in Polynomialzeit  lösbar.

Die Menge an Eingabefolgen, die einen Automaten synchronisieren, das heißt unser Suchraum, 
ist durch einen weiteren Automaton akzeptierbar, was die Untersuchung von akzeptierenden Automaton für jene Mengen nahelegt.
Von besonderem Interesse in Zusammenhang mit unseren Fragestellungen
ist der Fall, wenn ein akzeptierender Automat zur Beschreibung synchronisierender Eingabefolgen
maximal groß ist.
Wie wir hier herausarbeiten, wird die Größe solcher Automaten
durch zwei Faktoren beeinflusst: Erreichbarkeit und Unterscheidbarkeit.
Wir untersuchen beide Begriffe, stellen Zusammenhänge her
und setzen diese in Beziehung zu etablierten Kon\-zepten aus der Literatur, 
wie zum Beispiel zu
vollständig erreichbare Automaten (engl.~completely reachable automata).
Verschiedene Bedingungen für die Maximalität von Automaten für synchronisierende Eingabefolden
werden herausgearbeitet.
Weiterhin formulieren und untersuchen wir Entscheidungsprobleme und algorithmische Fragestellungen, 
welche sich aus diesen Begriffen ergeben.
Über ihren Transformationsmonoid haben Automaten einen natürlichen Bezug
zur Algebra.
Auf dieser Grundlage werden wir Verbindungen zur Theorie der Permutationsgruppen
herstellen. Im Speziellen
arbeiten wir, unter Zuhilfenahme graphentheoretischer Konstruktionen,
 verschiedene neuartige Charakterisierungen für primitive Permutationsgruppen heraus.
Unsere Ergebnisse verbinden Permutationsgruppen und ecken-transitive Graphen (engl.~vertex-transitive graphs) 
mit der Theorie synchronisierbarer Automaten.
Abschließend untersuchen wir, inwiefern die Größe eines akzeptierenden Automaten für die synchronisierenden Eingabefolgen
eines gegebenen Automaten mit der Länge einer kürzesten synchronisierenden Eingabefolge in Beziehung steht.

\selectlanguage{english}